\scanpage{172}
% L1–2
Considérons une extension
\diagram{\begin{tikzcd}[ampersand replacement=\&,column sep=small]
 \unsure{\cdot|0}\arrow[r,dash]\&E^0\arrow[r,dash]\&E\arrow[r,"\epsilon"]\&
 \underset{\substack{\rotatebox{-90}{$\simeq$}\\ E/E^0}}{F}\arrow[r,dash]\&0
\end{tikzcd}}
% L3–10
nous allons déterminer les involutions
de $E$, induisant $-\mathrm{id}_{E^0}$ sur $\unsure{E^0|E_0^0}$, et
$+\mathrm{id}_F$ dans $F$. Un endom \inserted{$\sigma$} de $E$ ayant
cette propriété (involutif ou pas) doit
être tel que
\[
u(x)=x+\sigma(x)=0\quad\text{si }x\in E^0,
\]
\margin{(exprime que $\sigma\mid E^0=-\mathrm{id}_{E^0}$}
donc $u$ se factorise en
\[
\pi:F\longrightarrow E\qquad u=\pi\cdot\epsilon
\]
% L11–12
et \deleted{$u(x)$} $\sigma$ se récupère via $\pi$ comme
\[
\deleted{\unsure{u(x|u(x)}} \sigma(x)=\pi(\epsilon(x))-x
\]

% L13–19
Si on exprime maintenant que
$\sigma$ induit $\mathrm{id}_F$ dans $E/E^0$, on trouve
la condition \(\displaystyle x\equiv\pi\epsilon(x))-x\quad(E^0)\quad\forall \unsure{1|x}\in E\)
i.e. \(\displaystyle \pi\underbrace{\epsilon(x)}_{y}\equiv2x\qquad\forall x\in E\)
ce qui signifie aussi que \unsure{le|les|la} $\epsilon$ des deux
\unsure{mbres|membres} est le $=$, soit (posant $\deleted{\unsure{\pi}}\epsilon(x)=y$
\[
\boxed{\epsilon\pi(y)=2y}
\]

% L20–22
Ainsi, $\pi$ est ``presque'' une section
c'est une $2$-section. \deleted{Cette} Plus \unsure{généralt|généralement},
si $\lambda\in k$ et si on a une application
% archive sheet number: 171
